\chapter{Bản chất công nghệ --- Bắt buộc lõi}
\label{ch:cong-nghe-cot-loi}

Chương này đi qua từng giai đoạn của chuỗi RAG, đối chiếu với đối tượng tương ứng của LlamaIndex,
và ghi lại kết quả chạy thật. Toàn bộ số liệu trong chương được sinh lại bằng
\texttt{make test} và \texttt{make eval} trên đúng phiên bản thư viện đã ghim trong
\texttt{requirements.txt} (xem Chương 3).


\section{Chuẩn bị dữ liệu: PDF scan thành Markdown hai tầng}
\label{sec:ingestion}

\subsection{Cơ chế}

LlamaIndex có sẵn \texttt{SimpleDirectoryReader} để đọc PDF, nhưng các văn bản quy chế của SGU là
bản scan: trang PDF chỉ là ảnh, không có lớp văn bản nên bộ đọc thường trả về gần như rỗng. Vì vậy
đồ án chuẩn bị dữ liệu ngoài LlamaIndex rồi mới đưa vào, theo hai tầng:

\begin{enumerate}
  \item \textbf{Tầng 1 --- \texttt{scripts/extract\_to\_md.py}}: nếu PDF có lớp văn bản thì trích trực tiếp;
        nếu không thì render từng trang ở 300\,dpi bằng \texttt{pdf2image} rồi OCR bằng
        \texttt{pytesseract} với \texttt{lang="vie"}~\cite{tesseract}. Script lưu kết quả vào
        \texttt{data/ocr/}, kèm chỉ số chất lượng là \emph{tỷ lệ âm tiết tiếng Việt hợp lệ}; khi một
        văn bản có nhiều bản scan, chỉ số này dùng để chọn bản tốt hơn.
  \item \textbf{Tầng 2 --- \texttt{scripts/clean\_md.py}}: làm sạch văn bản tầng 1 và dựng lại cấu trúc
        bằng heading Markdown: \texttt{\#} cho tên văn bản, \texttt{\#\#} cho Chương,
        \texttt{\#\#\#} cho Điều, \texttt{\#\#\#\#} cho mục sửa đổi. Kết quả nằm ở \texttt{data/processed/}.
\end{enumerate}

Tầng 1 là nơi \emph{duy nhất} được sửa tay; tầng 2 luôn sinh lại được từ tầng 1, nên không có
hai bản chân lý mâu thuẫn nhau. Cách tách này cũng là điều kiện để dùng lại toàn bộ hạ tầng
LlamaIndex phía sau: mọi thứ từ node trở đi đều đi trên Markdown có cấu trúc.

\subsection{Kết quả chạy thật}

Corpus hiện tại gồm 6 văn bản đang áp dụng. Chạy \texttt{python src/processed\_loader.py} cho:

\begin{lstlisting}[caption={Số Node theo từng văn bản (sinh lại được bằng \texttt{make test}).}, label={lst:node-theo-van-ban}]
54 Node từ 6 văn bản trong data/processed/01_quy_che_dao_tao/1_quy_che_dao_tao_dai_hoc
   5  2022-2375-QyD Chuyen Chuong Trinh Dao Tao DoiVoiSinhVenDHCQ.md
   1  22_quy-dinh-ve-khoi-luong-dao-tao-cua-cac-he-dao-tao-chinh-quy-dao-tao.md
  25  25_quy-che-dao-tao-trinh-do-dai-hoc-tai-truong-dai-hoc-sai-gon.md
  13  8_QD810QuyDinhVeHocCungLucHaiChuongTrinhDoiVoiSVDHCQ.md
   9  QD 3183.2025 - Sua doi Quy che Dao tao trinh do DH_LienQuanRutMonHoc.md
   1  Quydinh150tc.md
\end{lstlisting}

Văn bản dài nhất (Quy chế 2021) cho 25 Node: 21 Điều của Quy chế cộng 4 Node thuộc phần
``Quyết định ban hành''. Hai quy định một trang không chia Điều nên mỗi văn bản chỉ có một Node,
đúng như thiết kế \texttt{MIN\_TU\_NOI\_DUNG} (xem Mục~\ref{sec:node}).

\subsection{Vị trí trong đồ án và bài học từ bản cũ}

Giai đoạn 1 của bản báo cáo trước dùng \texttt{src/ingestion.py}: OCR từng trang và đóng gói
\emph{mỗi trang thành một \texttt{Document}}. Cách đó có hai nhược điểm đã đo được: (1) một Điều
nằm vắt qua hai trang bị chia đôi; (2) phần ``Quyết định ban hành'' ở trang đầu sinh ra các Node
trùng số Điều với phần Quy chế (Điều 1--3 xuất hiện hai lần), làm hỏng phép chấm truy hồi theo số
Điều. Hai module \texttt{src/ingestion.py} và \texttt{src/node\_parser.py} được giữ lại trong
repository để đối chiếu nhưng \textbf{không còn nằm trong pipeline}; dữ liệu vào Node ngày nay đi
qua tầng 2 nêu trên.

\section{Cắt đoạn theo cấu trúc văn bản (Node)}
\label{sec:node}

\subsection{Cơ chế}

Module \texttt{src/processed\_loader.py} đọc Markdown tầng 2 và sinh \texttt{TextNode} --- đối tượng
đơn vị của LlamaIndex~\cite{llamaindex2025docs}:

\begin{enumerate}
  \item Tách file theo heading bằng biểu thức chính quy \texttt{(?m)\^{}(\#\{1,4\}) } (lookahead giữ
        lại dòng heading), rồi theo dõi cấp heading để biết Node thuộc văn bản nào, Chương nào.
  \item Mỗi \texttt{\#\#\#} (Điều) hoặc \texttt{\#\#\#\#} (mục sửa đổi) là một Node; nhãn
        \texttt{dieu} hoặc \texttt{sua\_doi\_dieu} được trích bằng \texttt{Điều (\textbackslash d+)}.
  \item Nội dung nằm ngay dưới \texttt{\#}/\texttt{\#\#} chỉ thành Node khi dài từ
        \texttt{MIN\_TU\_NOI\_DUNG}~=~30 từ trở lên (dành cho văn bản không chia Điều).
  \item Node dài hơn \texttt{MAX\_TU}~=~1\,024 từ được cắt tiếp bằng \texttt{SentenceSplitter}
        (\texttt{chunk\_overlap}=50) và gắn thêm \texttt{phan\_doan}.
\end{enumerate}

\subsection{Hai sửa đổi quan trọng so với bản trước}

\paragraph{Loại metadata rỗng.} Bản trước gắn cho \emph{mọi} Node bốn khóa \texttt{van\_ban},
\texttt{chuong}, \texttt{dieu}, \texttt{sua\_doi\_dieu}, kể cả khi giá trị rỗng. Vì LlamaIndex ghép
metadata vào văn bản trước khi embedding, mỗi vector đều chứa các dòng vô nghĩa kiểu
\texttt{chuong:}, \texttt{dieu:} trong khi kho chỉ có vài chục Node. Bản mới lọc bỏ mọi khóa rỗng.

\paragraph{Chỉ một trường được embedding.} Bản mới sinh thêm khóa \texttt{nguon\_trich} và chỉ trường
này được đưa vào embedding; các khóa kỹ thuật còn lại bị loại bằng
\texttt{excluded\_embed\_metadata\_keys}. Văn bản đưa vào bộ mã hoá vì thế có dạng:

\begin{lstlisting}[caption={Văn bản đưa vào embedding của một Node (in bằng \texttt{metadata\_mode="embed"}).}, label={lst:embed-text}]
nguon_trich: Điều 9 — Quy chế đào tạo trình độ đại học SGU 2021

<nguyên văn Điều 9>
\end{lstlisting}

Trước đây cùng Node này cho chuỗi \texttt{van\_ban: ...\textbackslash nchuong: \textbackslash ndieu: 9\textbackslash nsua\_doi\_dieu:} ---
vừa nhiễu vừa chứa khóa tiếng Anh không dấu. Chỉ số \texttt{nguon\_trich} cũng chính là nhãn
\texttt{node\_label()} dùng để in nguồn, nên phần hiển thị và phần embedding không còn lệch nhau.
Đây là ví dụ cho thấy một quyết định nhỏ ở tầng Node ảnh hưởng trực tiếp tới chất lượng vector.

\subsection{Mã nguồn}

\begin{lstlisting}[language=Python, caption={Gắn metadata đã lọc rỗng và chọn trường đưa vào embedding (\texttt{src/processed\_loader.py}).}, label={lst:gan-metadata}]
def _gan_metadata(node: TextNode, meta: dict[str, str | int]) -> TextNode:
    """Gắn metadata đã lọc rỗng + `nguon_trich`, và đặt đúng trường nào được embed / vào prompt."""
    sach = {k: v for k, v in meta.items() if str(v).strip()}
    sach["nguon_trich"] = nguon_trich(sach)
    node.metadata = sach
    node.excluded_embed_metadata_keys = [k for k in sach if k != "nguon_trich"]
    node.excluded_llm_metadata_keys = [k for k in KHONG_LLM if k in sach]
    return node
\end{lstlisting}

\subsection{Kiểm thử}

\texttt{tests/test\_processed\_loader.py} khoá lại đúng những hành vi trên: số Node và nhãn Điều,
việc loại metadata rỗng, việc chỉ \texttt{nguon\_trich} được embedding, và giá trị của
\texttt{nguon\_trich} cho ba trường hợp (Điều thường, mục sửa đổi, phần ban hành).
\texttt{tests/test\_corpus.py} kiểm tra trên corpus thật: mọi tiền tố trong danh sách văn bản khớp
đúng một file, và mọi Node đều có \texttt{nguon\_trich}. Chạy \texttt{make test}: 26 kiểm thử, không
cần GPU và không cần Ollama.

\subsection{Hạn chế đã biết}

\begin{itemize}
  \item Cách cắt phụ thuộc \texttt{clean\_md.py} dựng heading đúng. Nếu một văn bản bị OCR làm mất
        dòng ``Điều $N$.'' thì Điều đó bị gộp vào Điều trước hoặc vào Node Chương.
  \item ``Token'' vẫn đếm bằng số từ tách theo khoảng trắng, chỉ là xấp xỉ.
  \item Hai văn bản một trang cho mỗi văn bản đúng một Node dài, chưa cắt theo khoản.
\end{itemize}

\section{Lập chỉ mục (Index)}
\label{sec:index}

\subsection{Cơ chế}

Module \texttt{src/index\_builder.py} làm ba việc:

\begin{enumerate}
  \item \textbf{Nạp Node} bằng \texttt{load\_corpus\_nodes()} của \texttt{processed\_loader}.
  \item \textbf{Đặt mô hình embedding} vào \texttt{Settings.embed\_model} --- \texttt{Settings} là
        cấu hình toàn cục của LlamaIndex, mọi thành phần cần embedding đều lấy từ đây. Hàm
        \texttt{resolve\_device()} tự hạ cấp về CPU khi máy không có CUDA, nhờ đó cùng một mã nguồn
        chạy được trên máy GPU, trên CPU và trên CI.
  \item \textbf{Dựng hoặc nạp lại} \texttt{VectorStoreIndex}. Chỉ mục được dựng thẳng từ danh sách
        \texttt{TextNode} đã cắt sẵn chứ không từ \texttt{Document}: nếu truyền \texttt{Document},
        LlamaIndex sẽ tự cắt bằng \texttt{SentenceSplitter} mặc định và làm mất ranh giới Điều ---
        đây chính là lý do phải gọi constructor bằng danh sách Node.
\end{enumerate}

\subsection{Fingerprint: chống dùng lại chỉ mục cũ}

Bản trước chỉ lưu tên mô hình embedding và sự tồn tại của \texttt{docstore.json}, nên sửa cách cắt
Node mà quên \texttt{--rebuild} sẽ âm thầm dùng chỉ mục cũ. Bản mới lưu \texttt{corpus.json} cạnh
chỉ mục, gồm: model embedding, thư mục corpus, \textbf{phiên bản loader cùng các ngưỡng cắt}, và
sha256 + tên hiển thị của từng file. Chỉ cần một trong các thứ đó đổi là chỉ mục tự build lại.

\begin{lstlisting}[language=Python, caption={Fingerprint của chỉ mục (\texttt{src/index\_builder.py}).}, label={lst:fingerprint}]
    return {
        "embed_model": cfg.embed_model_name,
        "corpus_dir": cfg.corpus_dir,
        "loader": {"phien_ban": LOADER_VERSION, "max_tu": MAX_TU, "min_tu_noi_dung": MIN_TU_NOI_DUNG},
        "files": files,
    }
\end{lstlisting}

\subsection{Kết quả chạy thật}

Với profile \texttt{server} (\texttt{AITeamVN/Vietnamese\_Embedding}, 1\,024 chiều) trên corpus 6
văn bản, lần build đầu ghi ra:

\begin{lstlisting}[caption={Dòng log khi build lại chỉ mục (đo trên CPU).}, label={lst:index-output}]
Corpus, model embedding hoặc loader đã đổi so với index đã lưu — build lại
Generating embeddings: 100%|##########| 54/54 [01:17<00:00, 1.43s/it]
Index: storage/server/01_quy_che_dao_tao/1_quy_che_dao_tao_dai_hoc | 54 Node
\end{lstlisting}

Chỉ mục lưu theo \texttt{storage/<profile>/<đường dẫn corpus>}, nên hai profile có số chiều vector
khác nhau không còn ghi đè lên nhau. Thời gian embedding 54 Node trên CPU khoảng 77 giây; trên GPU
4\,GB (đo trước đây) nhanh hơn khoảng một bậc.

\section{Truy hồi, trộn nguồn và xếp hạng lại}
\label{sec:retriever}

\subsection{Ba tầng của bước truy hồi}

Module \texttt{src/retriever.py} tổ chức truy hồi thành ba tầng độc lập, đúng theo cách LlamaIndex
tách \emph{retriever} khỏi \emph{node postprocessor}:

\begin{enumerate}
  \item \textbf{Vector top-$k$ (mặc định)}: \texttt{index.as\_retriever(similarity\_top\_k=10)}.
        Câu hỏi được embedding bằng đúng mô hình đã dùng khi lập chỉ mục, rồi so cosine với toàn bộ
        Node (\texttt{SimpleVectorStore} tìm vét cạn, phù hợp với kho vài chục đến vài trăm Node).
  \item \textbf{Hybrid BM25 + vector (tuỳ chọn, \texttt{use\_hybrid})}: câu hỏi về quy chế chứa nhiều
        ký hiệu và con số mà embedding nắm kém (``B+'', ``Điều 10'', ``150 tín chỉ'', ``QĐ 3183'').
        Đồ án cài BM25 Okapi~\cite{robertson2009bm25} thuần Python trong \texttt{src/bm25.py} ---
        gói \texttt{BM25Retriever} chính thức nằm ở một package riêng chứ không có trong
        \texttt{llama-index-core}, nên tự viết vừa nhẹ vừa cho phép kiểm soát cách tách từ.
        Hai nguồn được trộn bằng \texttt{QueryFusionRetriever} ở chế độ \texttt{reciprocal\_rerank}
        (RRF), với \texttt{num\_queries=1} để không gọi LLM sinh truy vấn phụ và giữ phép đo tất định.
  \item \textbf{Postprocessor}: ngưỡng tương đồng \texttt{SimilarityPostprocessor} (tuỳ chọn) rồi
        tới xếp hạng lại bằng cross-encoder~\cite{nogueira2019passage}
        (\texttt{SentenceTransformerRerank}, giữ top-3).
\end{enumerate}

\paragraph{Tách từ cho BM25.} \texttt{tach\_tu()} bỏ dấu tiếng Việt trước khi chấm điểm
(\texttt{Điều} $\rightarrow$ \texttt{dieu}). Nhờ vậy câu hỏi gõ thiếu dấu hoặc lỗi OCR vẫn khớp văn
bản có dấu --- đúng loại lỗi đã gặp trong nhật ký hỏi đáp. Đánh đổi là mất phân biệt các cặp chỉ
khác dấu, hiếm gặp trong văn bản quy chế. \texttt{tests/test\_bm25.py} khoá lại hành vi này bằng ba
ca: khớp theo ký hiệu ``B+'', khớp khi câu hỏi không dấu, và trả về rỗng khi không có từ nào khớp
(tránh đưa Node nhiễu vào prompt).

\paragraph{Một chi tiết dễ sai.} Điểm sau khi trộn RRF không cùng thang với cosine, nên nếu đặt
\texttt{similarity\_cutoff} khi đang bật hybrid thì mọi Node đều bị loại. Hàm
\texttt{get\_postprocessors()} vì thế chủ động bỏ qua ngưỡng khi \texttt{use\_hybrid=True} và ghi log
cảnh báo, thay vì trả về danh sách rỗng một cách im lặng.

\subsection{Kết quả chạy thật}

Trên bộ 31 câu hỏi về Quy chế 2021 (profile \texttt{server}, embedding trên CPU), cấu hình vector
cho \texttt{R@1 = 81\%}, \texttt{R@3 = 100\%}, \texttt{MRR = 0.892}. Phân bố văn bản của Node top-1:

\begin{lstlisting}[caption={Văn bản của Node đứng đầu, bộ 01, cấu hình vector.}, label={lst:top1-van-ban}]
Quy chế đào tạo trình độ đại học SGU 2021 : 23 câu
QĐ 3183/2025 sửa đổi Quy chế đào tạo SGU :  4 câu
QĐ 810 về học cùng lúc hai chương trình   :  3 câu
QĐ 2375/2022 về chuyển chương trình       :  1 câu
\end{lstlisting}

Bốn câu còn lại cho thấy tác dụng phụ của việc mở rộng corpus: văn bản khác chủ đề nhưng dùng từ
ngữ tương tự chen lên đầu. Trên bộ 02 (10 câu về nội dung sửa đổi) cả 10 câu đều lấy đúng QĐ 3183
ở top-1, nhưng câu hỏi về \emph{khóa tuyển sinh áp dụng} thang điểm mới trượt khỏi top-3 vì đáp án
nằm ở Điều 2 (điều khoản hiệu lực) trong khi truy hồi trả về các mục sửa đổi Điều 9, 10 --- chúng
khớp chữ ``B+'' nhiều hơn.

Kết quả đầy đủ của bốn cấu hình (vector, +rerank, hybrid, hybrid+rerank) và hai baseline nằm ở
Chương~\ref{ch:danh-gia}.

\section{Tổng hợp câu trả lời (Query Engine)}

\subsection{Cơ chế}

Module \texttt{src/query\_engine.py} ghép retriever, postprocessor và LLM thành
\texttt{RetrieverQueryEngine}:

\begin{itemize}
  \item \textbf{LLM}: \texttt{Ollama(model=..., thinking=False, context\_window=8192)}~\cite{ollama,qwen3}.
        Tắt chế độ suy nghĩ của Qwen3 để câu trả lời ngắn.
  \item \textbf{Tất định}: \texttt{temperature=0} và \texttt{seed=42} truyền xuống Ollama. Đây là
        sửa lỗi đã đo được ở bản trước: cùng một câu hỏi, lần đầu hệ thống từ chối, lần sau trả lời
        ``Không'', vì mô hình lấy mẫu ngẫu nhiên theo mặc định.
  \item \textbf{Prompt tiếng Việt riêng} thay prompt mặc định: nêu rõ ngữ cảnh là trích đoạn quy chế
        SGU, yêu cầu \emph{chỉ} dựa vào trích đoạn, và bổ sung ba ràng buộc nhắm đúng các lỗi đã ghi
        trong nhật ký: (1) \emph{trích nguyên văn câu quy định trước khi kết luận}, (2) \emph{không
        tự quy đổi thang điểm hay đổi đơn vị, cần tra bảng thì chép đúng hàng của bảng},
        (3) nói ``không tìm thấy'' khi trích đoạn không có thông tin --- cho mô hình một lối ra hợp
        lệ thay vì phải đoán.
  \item \textbf{Trích dẫn nguồn}: \texttt{format\_sources()} đọc metadata của các Node nguồn và in ra
        ``Nguồn: Điều X, Điều Y''. Nguồn lấy từ metadata chứ không để LLM tự ghi, nên luôn khớp với
        những gì đã đưa vào prompt.
  \item \textbf{Streaming}: câu trả lời in dần từng token qua \texttt{response.response\_gen}.
\end{itemize}

\subsection{Hạn chế của cấu hình này}

\texttt{RetrieverQueryEngine} chỉ gọi LLM \emph{một lượt} trên ngữ cảnh đã truy hồi: nó không nhớ
lịch sử hội thoại, nên câu hỏi nối tiếp kiểu ``còn kỹ sư?'' không đủ thông tin để truy hồi. Muốn
sửa phải đổi sang chat engine có viết lại câu hỏi theo lịch sử (ví dụ
\texttt{CondensePlusContextChatEngine} của LlamaIndex) --- việc này được nêu ở Chương~\ref{ch:danh-gia}
như hướng phát triển, và đã có sẵn lịch sử chat trong giao diện Gradio để dùng.

\section{Điểm ma sát với LlamaIndex quan sát được}
\label{sec:ma-sat}

Phần này ghi lại những chỗ thư viện không làm đúng như kỳ vọng ban đầu. Đây là loại thông tin
mà tài liệu chính thức không nêu, và là một trong những kết quả của việc \emph{dùng thật} thư viện.

\begin{table}[H]
  \centering
  \small
  \begin{tabularx}{\textwidth}{>{\hsize=1.05\hsize}X >{\hsize=1.05\hsize}X >{\hsize=0.9\hsize}X}
    \toprule
    \textbf{Hiện tượng} & \textbf{Nguyên nhân} & \textbf{Cách xử lý} \\
    \midrule
    Không nạp được embedding và reranker cùng lúc trên GPU 4\,GB &
      \texttt{SentenceTransformerRerank} không có tham số \texttt{dtype}, luôn nạp fp32 &
      Nạp reranker trên CPU, ép \texttt{\_model.half()} rồi mới chuyển lên GPU (dùng API nội bộ) \\
    Vector bị nhiễu bởi các trường metadata rỗng &
      Mặc định \texttt{metadata\_template} là \texttt{\{key\}: \{value\}} cho \emph{mọi} khóa, kể cả rỗng &
      Lọc khóa rỗng và chỉ để một trường (\texttt{nguon\_trich}) được embedding \\
    Không gắn được nhãn tiếng Việt cho từng khóa metadata &
      \texttt{metadata\_template} chỉ có hai chỗ điền \texttt{\{key\}}, \texttt{\{value\}}, không tùy biến theo khóa &
      Gom nhãn vào một trường \texttt{nguon\_trich} tự sinh \\
    Bật hybrid là lỗi \texttt{ImportError: llama-index-llms-openai} &
      \texttt{QueryFusionRetriever} làm \texttt{llm if llm else Settings.llm}, mà \texttt{Settings.llm} mặc định trỏ tới OpenAI &
      Truyền thẳng đối tượng \texttt{Ollama} vào retriever (với \texttt{num\_queries=1} nó không bao giờ được gọi) \\
    \texttt{BM25Retriever} không có trong \texttt{llama-index-core} &
      Thư viện tách phần này ra package riêng (\texttt{llama-index-retrievers-bm25}) &
      Tự viết \texttt{src/bm25.py} trên \texttt{BaseRetriever} (xem Mục~\ref{sec:retriever}) \\
    Không lọc được theo toán tử ``khác'' trên metadata &
      \texttt{SimpleVectorStore} mặc định chỉ là JSON trong RAM, hỗ trợ lọc hạn chế &
      Chấp nhận ở quy mô hiện tại; ghi nhận việc chuyển sang Chroma/Qdrant khi corpus lớn \\
    Đổi mô hình embedding là phải build lại toàn bộ chỉ mục &
      Định dạng lưu \texttt{default\_\_vector\_store.json} gắn với số chiều vector &
      Tách thư mục chỉ mục theo profile và thêm fingerprint theo model \\
    Cấu hình toàn cục gây tác dụng phụ &
      \texttt{Settings} là trạng thái global (embed\_model, llm) &
      Chỉ gán \texttt{Settings} ở đúng hai chỗ (\texttt{index\_builder}, \texttt{query\_engine}) \\
    \bottomrule
  \end{tabularx}
  \caption{Các điểm ma sát với LlamaIndex gặp trong quá trình làm đồ án.}
  \label{tab:ma-sat}
\end{table}

Một quan sát tổng quát: LlamaIndex tách bạch rất tốt ở mức \emph{interface} (retriever, node
postprocessor, response synthesizer đều thay được độc lập), nhưng ở mức \emph{mặc định} thì thiên
về ``chạy được ngay'' hơn là ``đúng cho tiếng Việt và cho văn bản pháp quy''. Vì vậy phần lớn công
sức của đồ án nằm ở việc thay các mặc định đó, chứ không ở việc ghép các thành phần lại.

\section{Tổng kết}

Sáu giai đoạn của pipeline đều đã có mã chạy được và có kiểm thử: chuẩn bị dữ liệu hai tầng, cắt
54 Node theo cấu trúc văn bản với metadata sạch, lập chỉ mục 1\,024 chiều có fingerprint theo loader,
truy hồi ba tầng (vector, tuỳ chọn hybrid BM25, tuỳ chọn rerank), sinh câu trả lời tất định có trích
dẫn Điều, và đánh giá tự động. So với bản báo cáo trước, ba vấn đề đã được xử lý và \emph{đo lại}:
Điều bị chia đôi giữa hai trang, nhãn Điều trùng giữa Quyết định ban hành và Quy chế, và câu trả lời
không ổn định giữa các lần hỏi.
